\subsection{Benchmark model and toolchain}\label{sec:model}
All measurements use a single benchmark model so that speedups are comparable across
optimizations: a single USGS HUC-8 watershed (hydrologic unit 03100101; the Peace River, southwest Florida), a
\emph{regional, hyper-resolution} SWAT+ model generated by the SWATGenX platform. The basin is
delineated from the USGS National Hydrography Dataset Plus High Resolution (NHDPlus~HR) over a
$1/3$~arc-second ($\approx$10~m) DEM, with land cover from the 500~m National Land Cover Database
(NLCD), soils from gridded SSURGO (gSSURGO), and climate forcing from PRISM (precipitation and
temperature) and the National Solar Radiation Database (NSRDB). The model comprises
57{,}998 hydrologic response units (HRUs), 8{,}181 channels, 9{,}341 routing units
(subbasins) resolved into 9{,}342 landscape units, 326 reservoirs and lakes, and 170 aquifers.
The high HRU count follows from the large drainage area combined with a hyper-resolution
slope/land-cover/soil intersection---this is a single regional model, not a continental one,
though the SWATGenX platform provides model-building coverage across the conterminous United
States (CONUS). Unless stated otherwise, the timed workload is one simulated year, which exercises
the full daily land-and-routing step over the entire network. The engine is built with the Intel
oneAPI Fortran compiler (\texttt{ifx}) at \texttt{-O3 -ipo}; \texttt{gfortran} fails to produce a
working binary at this model size, so \texttt{ifx} is the reference toolchain. OpenMP is enabled
with \texttt{-fiopenmp}. We keep two binaries---the unmodified serial engine and the parallel
engine---so the parallel work can be checked against an exact reference; at one thread the
parallel engine executes the original serial code path (Section~\ref{sec:validation}).

\subsection{The acceleration loop}
The methodology is a disciplined iteration, applied until returns diminish:
\begin{enumerate}
  \item \textbf{Profile.} Intel VTune \texttt{hotspots} collection on a representative window
        (180 simulated days), reported per routine and per source line, to find where time is
        actually spent rather than where it is assumed to be spent.
  \item \textbf{Classify.} Assign each hotspot to a class---startup name-matching, per-object
        re-initialization, derived-type copies, serial critical path, or genuine hydrologic
        compute---because the fix differs by class and the classification is the transferable
        result.
  \item \textbf{Fix by class.} Apply the class-appropriate change (e.g.\ index a name lookup,
        scope an array reset, privatize shared scratch, restructure for parallelism).
  \item \textbf{Validate.} Re-run the correctness battery (Section~\ref{sec:validation})
        before accepting any change.
  \item \textbf{Re-measure.} Re-time and re-profile; the new dominant hotspot becomes the next
        target. Each accepted change is a single version-control commit whose message records
        its validation evidence.
\end{enumerate}

\subsection{The routing DAG and the level wavefront}\label{sec:wavefront}
SWAT+ advances one daily time step by walking every object---HRU, routing unit, channel,
reservoir, aquifer---in a single fixed sequence (the \emph{command order}). That sequence
encodes the river network's connectivity: an object appears after all objects that feed it.
Equivalently, the command order is a topological sort of a directed acyclic graph (DAG) whose
edges are the hydrologic routing connections.

A topological sort admits a \emph{level} (antichain) decomposition: level $\ell$ contains
every object all of whose upstream dependencies lie in levels $<\ell$. Objects within a level
are mutually independent and may be computed concurrently; levels must be processed in order.
We precompute this level index once per run from the command order and drive the daily step
as a level wavefront: for each level, an OpenMP parallel loop computes its objects; an
implicit barrier between levels preserves the upstream-before-downstream ordering. For the
benchmark network this yields 265 levels: level~1 is the entire land phase
($\approx$58{,}000 HRUs---an enormous wide level with abundant parallelism), and the levels
narrow as the channel network converges, reaching width~1 at the deep main-stem
reaches---the serial critical path that ultimately bounds the achievable speedup
(Figure~\ref{fig:levels}; Section~\ref{sec:discussion}).

\begin{figure}[t]
\centering
\includegraphics[width=0.78\textwidth]{Figure_1.png}
\caption{The command-order level decomposition for the benchmark river network (265 levels).
Level~1 contains all 67{,}616 objects with no upstream dependency---the 57{,}998 HRUs together with
the other source (leaf) nodes (point-source/recall objects and headwater units)---an enormously wide,
abundantly parallel level; the levels narrow monotonically as the channel network converges,
reaching width~1 at the deep main-stem reaches (104 successive width-1 levels). The length of this
width-1 tail, not the per-level synchronization, sets the Amdahl ceiling on single-run speedup
(Section~\ref{sec:discussion}).}
\label{fig:levels}
\end{figure}

Two engineering refinements proved necessary. First, forking a parallel region \emph{per
level} (265 times per day) made thread-management overhead dominate and caused an
\emph{8-thread regression}; we instead open one parallel region per simulated day and
drive the levels inside it, with the per-level barrier as the only synchronization. Second,
the long tail of width-1 (and other very narrow) levels paid one barrier each for almost no
parallel work; we fuse consecutive narrow levels into a single master-thread serial
block, so the deep main stem pays one barrier instead of one-per-level
(Figure~\ref{fig:wavefront}).

\begin{figure}[t]
\centering
\includegraphics[width=0.82\textwidth]{Figure_2.png}
\caption{The routing-DAG level wavefront. The daily command order is a topological sort of the
routing DAG; objects at the same level are independent and run concurrently under a single
per-day OpenMP parallel region, with an implicit barrier between levels enforcing
upstream-before-downstream order. Consecutive narrow (near-width-1) levels on the main stem
are fused into one serial block to avoid paying a barrier per level. At one thread the
wavefront is inactive and the original serial command order executes unchanged.}
\label{fig:wavefront}
\end{figure}

\subsection{Reentrancy: privatizing current-object state}\label{sec:reentrancy}
The obstacle to running objects concurrently is not the routing math---it is that the engine,
though written in modern Fortran, communicates through \emph{module-global ``current-object''
scratch} by design: a routine sets a
global index (``the HRU we are on'', ``the channel we are on'') and downstream routines read
it. Under threads, concurrent objects clobber that shared state between the set and the use.
Making the per-object body reentrant therefore required auditing every such variable and
giving each thread its own copy via OpenMP \texttt{threadprivate}. Examples that had to be
privatized include the current weather-station index (a shared read of \texttt{iwst} between
its set and use caused $\sim$0.2\% PET drift on a station-dependent subset of HRUs), the
channel dispatch indices and rating-curve/peak-rate scratch, and a large cluster of
plant/residue ``temporary storage'' arrays.

A Fortran-specific hazard complicates this: a local declared with an initializer is
\emph{implicitly} \texttt{SAVE}d, hence shared across calls and therefore across threads, so the
reflexive ``strip the initializer'' fix is correct only for write-before-read scratch; the few
accumulators that legitimately persist must be kept \texttt{SAVE}d \emph{and} marked
\texttt{threadprivate}. This subtlety, and the one-thread byte-identity check that caught it, are
detailed in Supplementary Section~S2.

\subsection{Targeted hotspot removal (overhead, not hydrology)}\label{sec:hotspots}
Profiling repeatedly surfaced \emph{non-hydrologic} hotspots, each removed without touching the
science: an $\mathcal{O}(N^2)$ string name-match at startup (replaced by an $\mathcal{O}(1)$ index),
a per-object whole-array re-zeroing in the stream-temperature routine (\texttt{ch\_temp}, which
re-zeroed the entire landscape-unit array once per channel per day---the single hottest line in the
engine), and redundant derived-type struct copies in inflow-hydrograph handling. Each was scoped or
indexed to the data actually used, leaving output bit-identical; Section~\ref{sec:groupA} reports
the per-rung effect.

\subsection{Correctness framework}\label{sec:validation}
We validate with two complementary criteria.

\paragraph{Byte-identity at one thread.} At one thread the wavefront is inactive and the
engine executes the original command order, so its daily output must match the unmodified
serial engine \emph{bit for bit}. We verify this with exact byte-level diffs of every
daily HRU and channel record on a one-year benchmark run. This check caught a small
($\sim$0.3\%) \emph{deterministic} residual introduced by the reentrancy refactor in two
routines (\texttt{ch\_watqual4}, \texttt{et\_pot}), where stripping \texttt{SAVE} had dropped
persistent scratch values; restoring them as \texttt{threadprivate} \texttt{SAVE} state made
one-thread output bit-identical again (Section~\ref{sec:reentrancy}).

\paragraph{Thread-count invariance as a race detector.} For $N>1$ we require repeated runs at
different thread counts (1-vs-4, 4-vs-4, 4-vs-8) to produce identical daily output. Because a
correct shared-memory program is deterministic regardless of how work is scheduled, any
run-to-run difference is, by definition, a real data race---and the specific output fields
that differ point directly at the unprotected state. We use a strict tripwire (any record off
by $>10^{-6}$ counts as a difference). This criterion is what drove the iterative
race-elimination in Section~\ref{sec:reentrancy} and is fully automatable.

\paragraph{Separating compiler noise from incorrectness.} The \texttt{-ipo} production build
reorders floating-point operations, so two algorithmically identical builds can differ at the
last bits. We therefore distinguish (i) a non-\texttt{ipo} build, which is bit-identical to
the serial source and proves the \emph{algorithm} is unchanged, from (ii) the \texttt{-ipo}
build, whose residual ($\sim$$5\times10^{-4}$ relative, $\approx$0.008\%) is compiler
floating-point reassociation, not a logic change. Exercising the full process chain exposed one
$N>1$ data race in the in-stream \emph{particulate} constituents (Section~\ref{sec:correctness}),
localized to the channel constituent-transport path and under active resolution; hydrology and
dissolved species are model-equivalent across thread counts. The stream-temperature column is a
separate pre-existing upstream out-of-bounds read, output-only in the non-gwflow path, preserved
deliberately to match the reference serial behavior.
